\documentclass[10pt]{article}
\usepackage[margin=20mm]{geometry}
\usepackage{amsmath,amssymb,fontspec,unicode-math,xurl,hyperref}
\IfFontExistsTF{DejaVu Serif}{\setmainfont{DejaVu Serif}}{\setmainfont{Latin Modern Roman}}
\setmathfont{latinmodern-math.otf}
\setlength{\parindent}{0pt}\setlength{\parskip}{6pt}\sloppy
\begin{document}
ELVIA Scientific Verification Report 0073
Position: Stop Treating ’AGI’ as the North-star Goal of AI Research
Paper I25-40183 | Supported outcome: Minor
Review date: 24 October 2026 | Human scientific decision pending\par
\section*{1. Identity, outcome and scope}
Reviewed version: Borhane Blili-Hamelin and 15 coauthors, ICML 2025, PMLR 267:81090-81117. The frozen publisher-linked PDF has 28 pages; its title/authors and SHA-256 were verified against the publisher record and acquisition receipt. Locations use 1-based PDF pages, equal to printed internal page numbers. [R1]\par
Outcome: Minor: one bounded source-to-claim overstatement in the historical image-recognition example, with scope clarifications and unverified generalizations retained below. The review examined the main argument (pp.1-9), Appendices A-B and Table 1 (pp.25-27), and the contribution statement (p.28). Selected references were checked; the long bibliography was not exhaustively source-audited. This is not a certification of every historical example or whole-paper correctness.\par
\section*{2. Claim structure and decisive qualifications}
The six traps in Section 2 are a diagnosis supporting a normative proposal for specificity, pluralism and inclusion. They are explicitly interrelated rather than mutually exclusive (p.2). Section 2.1 does not treat contestation itself as a defect: it says disagreements can reveal conflicting values (p.3). Section 2.4 targets inadequately justified goals selected by incentives, rather than claiming that every hardware-dependent advance is scientifically illegitimate.\par
Appendix B (p.25) expressly distinguishes influential AGI advocates from a majority of AI researchers, and describes its quotations as a small illustrative sample. Table 1 is also explicitly illustrative and includes related concepts, not just definitions that use the acronym AGI. Neither appendix claims representative sampling or a prevalence estimate. No missing statistical test was invented for this material.\par
Section 4.1 (p.9) acknowledges that improved AGI approaches could mitigate the traps while retaining AGI as a goal. Footnote 9 expressly makes the reasons concerning specificity and pluralism defeasible. Section 3, footnote 8 (p.8), also allows that a unifying goal could be useful in some circumstances. Thus a hypothetical beneficial AGI project would not by itself refute the paper’s qualified argument. The report does not reinterpret the recommendation as an impossibility theorem or as evidence that every AGI program is harmful.\par
Morris et al.’s original paper (Introduction p.1; principles and framework p.4) explicitly separates goals, predictions and risks, and distinguishes performance depth from task breadth. This supports the target’s representation of Morris as a serious alternative, including Table 1 and Section 4.1, rather than an ignored counterexample. [R2]\par
\section*{3. Technical and evidentiary checks}
Underspecification (Section 2.2, p.3): D’Amour et al. define a pipeline-level problem in which multiple predictors satisfy validation criteria equally but behave differently in deployment. Their Introduction (pp.2-3) says this can occur even when the pipeline is well aligned with application goals. The target discusses a broader goal-specification problem. These notions should be distinguished: clearer language alone is not a verified cure for pipeline underspecification. The target supplies no theorem or empirical claim that goal clarification eliminates that failure mode, so this distinction does not establish a false technical guarantee. [R3]\par
Generality debt (Section 2.5, p.5): the target lists different meanings of generality and says certain forms could trade off against ecological validity. It acknowledges that methodological challenges may be overcome. Saxon et al., Section 2.2 (pp.3-4), discusses precisely the gap between benchmark collections and application-relevant evaluation. Paolo et al., Section 5.5 (arXiv v4, pp.8-9), discusses specifying adaptation to unseen settings. These checks support the conceptual concern, not a universal no-generalization theorem. The target’s stronger prevalence wording about most papers lacks an independently checked quantitative denominator; no representative literature audit was performed. [R4-R5]\par
Compute growth (Section 2.6, pp.6-7): Sevilla et al., Table 2 (p.3), reports 5.7-month doubling for selected regular-scale models during 2010-2022 and 9.9 months for its large-scale group during late 2015-2022. The target’s broad wording about growth over months is consistent with rapid historical growth, but does not establish a minimum compute requirement for every AI project or a forecast. No regression or training-cost experiment was reproduced. [R6]\par
Specificity example (Section 3, p.7): the short/medium/long utterance example is explicitly illustrative. It contrasts a broad research question with a domain-specific question; it does not claim measured superiority for the hypothetical mixture of experts. Similarly, pluralism and inclusion recommendations are supported by arguments and selected studies, not a new controlled intervention establishing universal efficacy.\par
\section*{4. Problem 1 - Historical recurrence is stronger than the checked evidence}
The actual sentence in Section 2.6 (p.6) says that engines mistaking Black people for gorillas are still doing so more than eight years after the problem was identified. This is a claim of ongoing harmful outputs, not merely a statement that an image-recognition limitation persists.\par
The cached Appelman (2023) account describes the 2015 Google mislabeling and a 2021 Facebook primates incident. Its new 2023 observation, reporting a 44-image test, is that Google and Apple returned no gorilla search results; it attributes this to disabled primate search. That observation demonstrates a continuing restriction or workaround. It does not itself establish fresh harmful labels in 2023. The demonstrated defect is this source-to-claim overstatement. Actual 2023 recurrence remains Needs Evidence; the check does not prove that recurrence was false or that the underlying bias was fixed. The full primary test record was not independently acquired. [R7-R8]\par
Counterevidence and dependency: the charitable reading is persistence of an unresolved recognition problem, and the workaround supports that narrower concern. Earlier harmful incidents also remain documented. But those qualifications cannot silently replace the paper’s literal recurrence claim. The defect affects one historical illustration in the neglect and exclusion discussion. The six-trap argument, the specificity/pluralism/inclusion recommendations, and the explicitly defeasible discussion of alternatives do not require this particular temporal assertion. This bounded evidentiary overstatement supports Minor.\par
Scoped correction: distinguish the 2015/2021 incidents from the 2023 restricted-search observation, and describe the latter as a continuing workaround rather than a demonstrated fresh misclassification. No corrected model accuracy, current-world safety claim, or causal impact estimate is supplied.\par
Other broad empirical language, including the extent to which AGI discourse causes resource concentration or worsens each trap, was not independently quantified. Such questions would require a specified population, measures and causal design. This limits evidentiary certification; it is not an added experimental prerequisite for the normative recommendation.\par
\section*{5. Versions, code and conclusion}
The inspected arXiv record lists v1 on 6 February 2025, v3 on 6 May 2025 and v4 on 7 July 2025. The v4 HTML was consulted for provenance and selected scope language; it was not substituted for the frozen publisher PDF and byte equivalence is not asserted. No later erratum was identified in the inspected records. OpenReview returned a browser-verification challenge, so its discussion and revision history were not audited. [R9]\par
No original executable method, dataset analysis or author repository was identified in the inspected publisher page and primary text. No author code was run. The historical image example has a bounded Minor source-to-claim overstatement; actual 2023 recurrence remains Needs Evidence. No Major defect was established, and the other inspected arguments retain their scoped assessment. No author adoption is claimed.\par
Inspection | Outcome | Repair / adoption | Remaining work\par
Other inspected arguments and appendix scope | No additional issue | No correction proposed | Normative scientific decision pending\par
Other selected technical sources | No additional issue, qualified | No numerical result repaired | Full cited-study and bibliography audit not done\par
Underspecification | Scope clarification | No verified repair needed | Do not infer a universal cure for pipeline underspecification\par
Human review | Pending | Adoption not requested | publication status and human judgment\par
Historical image example | Minor; recurrence Needs Evidence | Separate incidents from 2023 suppression | Do not claim observed fresh labels or universal absence\par
\section*{6. Consulted references}
[R1] Blili-Hamelin et al. ICML 2025, PMLR 267:81090-81117. https://proceedings.mlr.press/v267/blili-hamelin25a.html .\par
[R2] Morris et al. Position: Levels of AGI for Operationalizing Progress on the Path to AGI. ICML 2024, PMLR 235:36308-36321. Publisher PDF, pp.1,4: https://proceedings.mlr.press/v235/morris24b.html .\par
[R3] D’Amour et al. Underspecification Presents Challenges for Credibility in Modern Machine Learning. JMLR 23(226):1-61, 2022. Introduction pp.2-3: https://jmlr.org/papers/v23/20-1335.html .\par
[R4] Saxon et al. Benchmarks as Microscopes: A Call for Model Metrology. COLM 2024 accepted manuscript, arXiv:2407.16711v2, 30 July 2024, Section 2.2. https://arxiv.org/abs/2407.16711 .\par
[R5] Paolo et al. A Call for Embodied AI. arXiv:2402.03824v4, 13 September 2024, Section 5.5. The exact cited conference version was not independently compared. https://arxiv.org/abs/2402.03824 .\par
[R6] Sevilla et al. Compute Trends Across Three Eras of Machine Learning. arXiv:2202.05924v2, 9 March 2022, Table 2. https://arxiv.org/abs/2202.05924 .\par
[R7] Grant and Hill. Google’s Photo App Still Can’t Find Gorillas. And Neither Can Apple’s. New York Times, 22 May 2023; original bylined reporting retrieved through search-indexed syndication: https://www.ekathimerini.com/nytimes/1212118/ . Direct opening returned 403. [R8] Appelman, 9 June 2023, cited commentary and links: https://racismandtechnology.center/2023/06/09/racist-technology-in-action-image-recognition-is-still-not-capable-of-differentiating-gorillas-from-black-people/ .\par
[R9] Version record: https://arxiv.org/abs/2502.03689 ; v4 HTML: https://arxiv.org/html/2502.03689v4 . Retrieved 24 October 2026.\par
\section*{Supporting checks}
Primary SHA-256: `daf5ecc93a7b36f4b8b793606cf81b1cb36bc7fbdda27cc4957f88ecbdfd5f6a`. Original bytes are preserved.\par
No empirical reproduction or author-code execution was performed.\par
Recorded failures: initial extraction could not create text because sources/0073 did not yet exist; after inspecting the error, the directory was created and extraction completed. A legacy Morris PDF URL returned 404; the publisher-linked raw URL succeeded. An inline helper failed at Python parsing because of shell quoting and was replaced by inspect\_sources.py. OpenReview and direct syndicated-news access were limited as described above.\par
No human verification is declared. Scientific decision is pending and author adoption has not been requested.\par
\section*{Supporting checks}
This proposed replacement was prepared during accept-0007. It changes only the outcome, the historical image-recognition assessment and the corresponding summary entry. The exact predecessor hashes for the six input artifacts are retained in `predecessor\_hashes.json`. All other scientific assessments and existing source limitations are preserved. The supporting files remain in the frozen original staging tree. No human scientific adjudication or archive promotion is implied.\par
\end{document}
